\section{Paso 6: validación del modelo programado}
\label{sec:paso6-validacion}
% Origen: 07-plan-tpi.md §6.2, 02-modelo-anylogic.md (piloto del 05/10 y reproducción del 09/10) y 09-pedidos-de-datos.md (respuesta de Trenes Argentinos, 26/10).

Validar es comprobar que el modelo representa el sistema real con la precisión que exige la pregunta
\parencite{law2015}. La verificación (sección~\ref{sec:paso4-construccion-verificacion}) comprueba otra
cosa: que el programa hace lo que el modelo especifica. Con los datos disponibles, el modelo puede
contrastarse con la realidad en el caudal, en el uso de los molinetes y en la llegada de los trenes. No
puede contrastarse en la espera ni en el largo de cola. La tabla~\ref{tab:validacion} resume cada control.

\begin{table}[H]
\centering
\caption{Controles de validación}
\label{tab:validacion}
\small
\begin{tabularx}{\textwidth}{@{}>{\raggedright\arraybackslash}p{3.3cm} X >{\raggedright\arraybackslash}p{2.6cm}@{}}
\toprule
Control & Contra qué se compara & Resultado \\
\midrule
Caudal por ventana & Validaciones simuladas frente al perfil de SBASE & Calibración, no validación (ver texto) \\
Molinetes en uso & Molinetes con tráfico en \EO{} frente a los 19,5 activos observados & Consistente \\
Reparto entre molinetes & Utilización por molinete frente al desbalance observado & Más parejo que la realidad \\
Llegadas de los trenes & Horario oficial frente a llegadas reales & Atraso no modelado \\
Volumen que llega en tren & Demanda del modelo frente a la matriz origen-destino del Roca & Verosímil \\
Geometría del andén & Dispersión de la descarga frente al plano ferroviario & Del mismo orden \\
Espera y cola & Sin observaciones publicadas & No validable \\
\bottomrule
\end{tabularx}
\end{table}

\subsection{Caudal por ventana}

El modelo cuenta las validaciones simuladas en cada ventana de 15 minutos y las compara con el perfil de
SBASE. Este control es en parte circular, porque la demanda se construye con ese mismo perfil. Por eso se
presenta como calibración: comprueba que la demanda se reparte bien entre los trenes y que la congestión no
retiene pasajeros de una ventana a la siguiente. Con valores de prueba, el caudal simulado quedó en su
mayoría a menos del 6~\% del observado (sección~\ref{sec:paso5-pilotos}).
\pendiente{registrar el caudal por ventana con el escenario intermedio (consola de PeatonalFranjaVisualE0).}

Una validación no circular sería correr el modelo con el perfil de un subperíodo que no se usó para
calibrarlo y comparar. Requiere que el perfil cambie entre réplicas, lo que el diseño actual no hace
(supuesto~4 de la sección~\ref{sec:paso2-datos-modelo-conceptual}).

\subsection{Uso de los molinetes}

\textbf{Molinetes en uso.} En el piloto de \EO{} se usaron los 20 molinetes habilitados. Coincide con los 19,5 molinetes
con tráfico por ventana que registra SBASE.

\textbf{Reparto.} En la realidad, el molinete más usado procesa 11 veces más pasajeros que el menos usado
\parencite{sbase2026molinetes}. En el piloto de \EO{}, las utilizaciones fueron de entre 17 y 52~\%, un
factor de alrededor de 3. \textcite{seriani2025} validan su modelo justamente con esta comparación. El
modelo reparte la carga más parejo que la realidad, y hay dos razones posibles:
\begin{itemize}
  \item la geometría es hipotética, y en la estación real algunos molinetes quedan lejos del recorrido
        natural;
  \item la regla de elección pondera la fila tanto como la distancia.
\end{itemize}
La consecuencia importa para la conclusión. Si los dos molinetes que agrega \EI{} quedaran, en la
estación real, en una posición poco atractiva, se usarían menos que en el modelo. En ese caso, el modelo
\emph{sobreestimaría} la mejora de \EI{}.

\subsection{Llegadas reales de los trenes}

Trenes Argentinos informó la llegada real de cada tren a Plaza Constitución entre el 02/03 y el 30/06/2026,
de 07:00 a 09:30 \parencite{sofse2026nota}. Es el mismo período del perfil de SBASE. En los 76 días
hábiles normales llegaron entre 43 y 52 trenes por día, con una media de 48,9. El modelo usa 47.
Comparado con el horario vigente, el atraso tiene estas características:
\begin{itemize}
  \item mediana de $+4$~minutos;
  \item percentil 10 de $-1$~minuto y percentil 90 de $+11$~minutos;
  \item percentil 99 de $+24$~minutos.
\end{itemize}
Además, el tren 302, de larga distancia desde Mar del Plata, llega todos los días en la franja y no figura
en los cuadros horarios usados.

Un atraso parejo solo desplaza las oleadas en el tiempo. Un atraso variable, en cambio, junta trenes que
el horario separa. Es esperable que eso produzca oleadas más grandes y más cola que las del modelo. El
supuesto~5 (trenes en horario) es entonces probablemente optimista sobre la espera. Este dato permite
reemplazar el horario por las llegadas reales de cada día, que es la extensión natural del modelo.

\subsection{Volumen que llega en tren}

La matriz origen-destino del Roca indica que entre el 44 y el 60~\% de los viajes de día hábil que salen
de cada estación de la línea terminan en Constitución. Aplicada a los pasajeros que suben en cada estación
entre las 07:00 y las 09:59, da unos \num{36000} pasajeros por día hábil que llegan a la terminal,
cerca de 700 por tren. En el escenario intermedio, el modelo hace entrar por el vestíbulo Principal
$0{,}80 \times \num{18532} \approx \num{14800}$ pasajeros desde los trenes. Haría falta que cerca del 40~\%
de los que llegan pasen a la Línea~C. La proporción es verosímil para la principal conexión de la
terminal con el centro, pero no la confirma ningún dato: Trenes Argentinos informó que no hay estudios de
trasbordo. El valor de 0,80 queda como supuesto, ahora con un orden de magnitud que lo respalda.

\subsection{Geometría del andén}

El plano de la estación enviado por Trenes Argentinos muestra los paragolpes de todas las vías alineados y
andenes útiles de entre 134 y 275~m. La dispersión de la descarga del escenario intermedio (145~s) supone
un recorrido de unos 138~m entre el primer y el último coche, del mismo orden que los andenes más cortos.
El plano no incluye los accesos ni el vestíbulo de la Línea~C, que pertenecen a SBASE, por lo que no sirve
para validar la geometría del vestíbulo hipotético.

\subsection{Reproducción en otra máquina}

Los 30 pares de producción se repitieron el 09/10/2026 en otra computadora, con Windows, el mismo archivo
del modelo y las mismas semillas. La demanda generada coincidió en las 60 corridas. El movimiento no fue
idéntico bit a bit: las esperas difirieron en décimas de segundo, probablemente por la aritmética de coma
flotante de la simulación peatonal. Las diferencias \EI{} $-$ \EO{} de las tres medidas primarias variaron
menos de 0,15~s, y la proporción coincidió con tres decimales. Los resultados no dependen de la máquina.

\subsection{Lo que no se puede validar}

No hay observaciones publicadas de la espera ni del largo de cola en Constitución. Los molinetes registran
a quien pasó, no a quien esperó. El grupo no pudo medir en el lugar, y el pedido a SBASE incluye
mediciones de cola si existen. Es la limitación principal del estudio. Los valores absolutos de espera
dependen de los supuestos de la tabla~\ref{tab:parametros} y de la geometría hipotética. La comparación
entre escenarios es más robusta que los valores absolutos, porque los dos escenarios comparten todos esos
supuestos. Aun así, el reparto más parejo del modelo puede inflar la diferencia.

En síntesis, el modelo reproduce lo que el dato permite observar: el caudal, la cantidad de molinetes en
uso y un volumen en tren compatible con la carga del Roca. Se aparta de la realidad en dos puntos
identificados: reparte la carga más parejo y supone trenes en horario. Las conclusiones de las secciones
siguientes se formulan como condicionales a esos supuestos.
